\documentclass[12pt, letterpaper]{article}

% Paquetes esenciales (coherentes con tu libro)
\usepackage[utf8]{inputenc}
\usepackage[spanish, es-nodecimaldot]{babel}
\usepackage{amsmath, amssymb, amsthm}
\usepackage{geometry}
\usepackage{xlop}
\usepackage{tikz}
\usepackage{cancel}
\usepackage{xcolor}

% Configuración de página
\geometry{left=2.5cm, right=2.5cm, top=2.5cm, bottom=2.5cm}

% Definición de entornos para soluciones
\newenvironment{solucion}
    {\begin{proof}[Solución]}
    {\end{proof}}

\title{Solucionario Detallado: Sistemas Numéricos y Teoría de Números}
\author{Matemáticas Básicas}
\date{}

\begin{document}

\maketitle

\section*{Nivel I: Mecánico y Procedimental}

\subsection*{1. Cálculo de $\mathcal{MCD}$ y $\mathcal{MCM}$}

\textbf{Ejercicio 1.a)} $a = 450, \quad b = 360$

\begin{solucion}
    \textbf{Paso 1: Descomposición en factores primos.}
    \begin{center}
    \begin{tabular}{c|c} 
        450 & 2 \\ 225 & 3 \\ 75 & 3 \\ 25 & 5 \\ 5 & 5 \\ 1 & 
    \end{tabular} \quad $\Rightarrow 450 = 2^1 \cdot 3^2 \cdot 5^2$
    \hspace{1cm}
    \begin{tabular}{c|c} 
        360 & 2 \\ 180 & 2 \\ 90 & 2 \\ 45 & 3 \\ 15 & 3 \\ 5 & 5 \\ 1 & 
    \end{tabular} \quad $\Rightarrow 360 = 2^3 \cdot 3^2 \cdot 5^1$
    \end{center}

    \textbf{Paso 2: Cálculo del $\mathcal{MCD}$ (factores comunes con menor exponente).}
    \[ \mathcal{MCD}(450, 360) = 2^1 \cdot 3^2 \cdot 5^1 = 2 \cdot 9 \cdot 5 = \mathbf{90} \]

    \textbf{Paso 3: Cálculo del $\mathcal{MCM}$ (comunes y no comunes con mayor exponente).}
    \[ \mathcal{MCM}(450, 360) = 2^3 \cdot 3^2 \cdot 5^2 = 8 \cdot 9 \cdot 25 = \mathbf{1800} \]
\end{solucion}

\textbf{Ejercicio 1.b)} $a = 792, \quad b = 900$

\begin{solucion}
    $792 = 2^3 \cdot 3^2 \cdot 11^1$ \\
    $900 = 2^2 \cdot 3^2 \cdot 5^2$
    
    \[ \mathcal{MCD} = 2^2 \cdot 3^2 = 4 \cdot 9 = \mathbf{36} \]
    \[ \mathcal{MCM} = 2^3 \cdot 3^2 \cdot 5^2 \cdot 11^1 = 8 \cdot 9 \cdot 25 \cdot 11 = \mathbf{19800} \]
\end{solucion}

\subsection*{2. Algoritmo de Euclides e Identidad de Bézout}

\textbf{Ejercicio 2.a)} $a = 123, \quad b = 277$ (Hallar $\mathcal{MCD}$ y expresar combinación lineal).

\begin{solucion}
    Aplicamos el algoritmo de Euclides dividiendo el mayor por el menor:
    \begin{align}
        277 &= 123(2) + 31 \label{eq:euc1} \\
        123 &= 31(3) + 30 \label{eq:euc2} \\
        31 &= 30(1) + 1 \label{eq:euc3} \\
        30 &= 1(30) + 0 \nonumber
    \end{align}
    El último residuo no nulo es $1$, por lo tanto $\mathcal{MCD}(123, 277) = \mathbf{1}$.

    \textbf{Reversa (Bézout):} Despejamos el residuo en cada paso.
    De (\ref{eq:euc3}): $1 = 31 - 30(1)$
    
    Sustituimos $30$ usando (\ref{eq:euc2}): $30 = 123 - 31(3)$
    \[ 1 = 31 - [123 - 31(3)] = 31(1) - 123 + 31(3) = 31(4) - 123 \]
    
    Sustituimos $31$ usando (\ref{eq:euc1}): $31 = 277 - 123(2)$
    \[ 1 = [277 - 123(2)](4) - 123 = 277(4) - 123(8) - 123(1) = 277(4) - 123(9) \]
    
    \textbf{Resultado:} $1 = 277(4) + 123(-9)$.
\end{solucion}

\section*{Nivel II: Análisis y Aplicaciones}

\subsection*{5. Problema de Baldosas}
Habitación de $420 \times 480$ cm.

\begin{solucion}
    Se busca el tamaño máximo cuadrado que divida exactamente a ambas dimensiones. Esto corresponde al $\mathcal{MCD}(420, 480)$.
    
    Descomposición rápida:
    $420 = 10 \cdot 42 = 2 \cdot 5 \cdot 2 \cdot 3 \cdot 7 = 2^2 \cdot 3 \cdot 5 \cdot 7$
    $480 = 10 \cdot 48 = 2 \cdot 5 \cdot 16 \cdot 3 = 2^5 \cdot 3 \cdot 5$
    
    \textbf{a) Tamaño de la baldosa:}
    $\mathcal{MCD} = 2^2 \cdot 3 \cdot 5 = 4 \cdot 3 \cdot 5 = \mathbf{60 \text{ cm}}$.
    
    \textbf{b) Cantidad de baldosas:}
    \[ \text{Largo: } \frac{420}{60} = 7 \text{ baldosas}, \quad \text{Ancho: } \frac{480}{60} = 8 \text{ baldosas} \]
    \[ \text{Total} = 7 \times 8 = \mathbf{56 \text{ baldosas}}. \]
\end{solucion}

\subsection*{6. Ecuación Diofántica (Granjero)}
Ecuación: $100x + 17y = 1800$, con $x, y \in \mathbb{Z}_{\ge 0}$.

\begin{solucion}
    Primero verificamos si tiene solución. $\mathcal{MCD}(100, 17) = 1$ (17 es primo y no divide a 100). Como $1 | 1800$, hay solución.
    
    Usamos congruencias (o despeje simple) para reducir el problema. Trabajamos módulo 17:
    \[ 100x + 17y \equiv 1800 \pmod{17} \]
    Como $17y \equiv 0 \pmod{17}$:
    \[ 100x \equiv 1800 \pmod{17} \]
    Reducimos coeficientes:
    $100 = 17(5) + 15 \Rightarrow 100 \equiv 15 \equiv -2 \pmod{17}$.
    $1800 = 17(105) + 15 \Rightarrow 1800 \equiv 15 \equiv -2 \pmod{17}$.
    
    La ecuación queda:
    \[ -2x \equiv -2 \pmod{17} \implies x \equiv 1 \pmod{17} \]
    
    Por lo tanto, $x$ tiene la forma $x = 17k + 1$.
    Como $x$ representa número de vacas, $x \ge 0$.
    Probemos $k=0 \Rightarrow x=1$.
    Sustituimos en la original: $100(1) + 17y = 1800 \Rightarrow 17y = 1700 \Rightarrow y = 100$.
    
    Si probamos $k=1 \Rightarrow x=18$.
    $100(18) + 17y = 1800 \Rightarrow 1800 + 17y = 1800 \Rightarrow y = 0$.
    
    \textbf{Soluciones posibles:}
    1. 1 vaca y 100 ovejas.
    2. 18 vacas y 0 ovejas.
\end{solucion}

\section*{Nivel III: Teórico y Demostraciones}

\subsection*{7. Propiedad $\mathcal{MCD}(ka, kb) = k \cdot \mathcal{MCD}(a, b)$}

\begin{solucion}
    Sea $d = \mathcal{MCD}(a, b)$ y $D = \mathcal{MCD}(ka, kb)$.
    
    1. Por definición de $\mathcal{MCD}$, podemos escribir $a = d \cdot x$ y $b = d \cdot y$, donde $\mathcal{MCD}(x, y) = 1$.
    
    2. Multiplicamos por $k$:
    $ka = k(dx) = (kd)x$
    $kb = k(dy) = (kd)y$
    
    3. Observamos que $kd$ es un divisor común de $ka$ y $kb$.
    Por lo tanto, $\mathcal{MCD}(ka, kb) = (kd) \cdot \mathcal{MCD}(x, y)$.
    
    4. Como $\mathcal{MCD}(x, y) = 1$, entonces:
    \[ \mathcal{MCD}(ka, kb) = kd \cdot 1 = k \cdot d = k \cdot \mathcal{MCD}(a, b) \]
\end{solucion}

\subsection*{10. El Algoritmo de la División ($n^2-1$ divisible por 8)}
Demostrar que si $n$ es impar, $8 | (n^2 - 1)$.

\begin{solucion}
    Si $n$ es impar, podemos escribirlo como $n = 2k + 1$ para algún $k \in \mathbb{Z}$.
    Sustituimos en la expresión:
    \[ n^2 - 1 = (2k+1)^2 - 1 = (4k^2 + 4k + 1) - 1 = 4k^2 + 4k \]
    Factorizamos $4k$:
    \[ n^2 - 1 = 4k(k+1) \]
    
    Analicemos el término $k(k+1)$. Es el producto de dos enteros consecutivos.
    En cualquier par de enteros consecutivos, uno de ellos debe ser par.
    Por lo tanto, $k(k+1)$ es divisible por 2, es decir, $k(k+1) = 2m$ para algún entero $m$.
    
    Sustituyendo esto atrás:
    \[ n^2 - 1 = 4(2m) = 8m \]
    
    Esto demuestra que $n^2 - 1$ es un múltiplo de 8. $\qed$
\end{solucion}

\section*{Nivel IV: Retos}

\subsection*{12. Sucesión de Fibonacci y $\mathcal{MCD}$}
Calcule $\mathcal{MCD}(F_{n}, F_{n+1})$.

\begin{solucion}
    Usamos el algoritmo de Euclides.
    Sabemos que $F_{n+1} = F_n \cdot 1 + F_{n-1}$ (por definición de la sucesión).
    
    Por la propiedad del algoritmo de Euclides: $\mathcal{MCD}(a, b) = \mathcal{MCD}(b, r)$ donde $a = bq + r$.
    
    Aplicando esto a Fibonacci:
    \[ \mathcal{MCD}(F_{n+1}, F_n) = \mathcal{MCD}(F_n, F_{n-1}) \]
    
    Podemos repetir este proceso iterativamente (descenso infinito):
    \[ \mathcal{MCD}(F_n, F_{n-1}) = \mathcal{MCD}(F_{n-1}, F_{n-2}) = \dots = \mathcal{MCD}(F_2, F_1) \]
    
    Sabemos que $F_1 = 1$ y $F_2 = 1$.
    \[ \mathcal{MCD}(1, 1) = 1 \]
    
    \textbf{Conclusión:} Dos números consecutivos de Fibonacci son siempre primos relativos ($\mathcal{MCD} = 1$).
\end{solucion}

\end{document}
